\documentclass[italian,12pt]{article}

\def\Title{Verbale interno}
\def\Subtitle{Riunione interna settimanale}
\def\Author{7Last}
\def\Date{2024-05-22}
\def\Version{v1.0}

%pacchetti extra da scaricare dblfloatfix, fancyhdr
\usepackage[left=2cm, right=2cm, bottom=3cm, top=3cm]{geometry}
\usepackage{fancyhdr}%creazione header-footer
\usepackage{graphicx} %serve per inserire immagini
\graphicspath{ {../../../logo/} }
%\usepackage{dblfloatfix} %serve per posizionare gli elementi dove si vuole
\usepackage[hidelinks]{hyperref} %serve per i link
\usepackage{tikz}
\usepackage{tgadventor} % font
\usepackage[useregional=numeric,showseconds=true,showzone=false]{datetime2}
\usepackage{caption}

\usepackage{hyperref}
\usepackage{tocloft}
\usepackage{titlesec}
\usepackage{color}
\usepackage{ulem}
\usepackage{pgfplots}
\usepackage{pgf-pie}
\usepackage[italian]{babel}
\usepackage{comment}
\usepackage{tabularx}
\usepackage{longtable}
\usepackage{float}
\usepackage{amsmath}

\linespread{1.2}
\captionsetup[table]{labelformat=empty}
\geometry{headsep=1.5cm}

\renewcommand{\contentsname}{Indice}
\renewcommand\familydefault{\sfdefault}

\let\oldthepage\thepage
\renewcommand{\thepage}{\sffamily \oldthepage}

\begin{document}

\newgeometry{left=2cm,right=2cm,bottom=2.1cm,top=2.1cm}
\begin{titlepage}
    \vspace*{.5cm}

    \vspace{2cm}
    {
        \centering
        {\bfseries\huge \Title\par}
        \bigbreak
        {\bfseries\large \Author\par}
        \bigbreak
        {\Version\par}
        \vfill

        \begin{tikzpicture}[remember picture,overlay]

            \fill[blue!20!black] (current page.south west) -- ++(10cm,0) -- ++(-10cm,15cm) -- cycle;
            \fill[orange] (current page.south east) -- ++(-18cm,0) -- ++(21.6cm,22cm) -- cycle;

            \clip (0,-3cm) circle (2.5cm) node (current page.center) {\includegraphics[width=5cm]{logo.jpg}};
        \end{tikzpicture}

    }

    \vfill

\end{titlepage}

\restoregeometry






















\newpage

%---------------header------------------%
\pagestyle{fancy}
\fancyhead{} % pulizia degli header
\lhead{%
\begin{tikzpicture}
    \clip (0,0) circle (0.5cm);
    \node at (0,0) {\includegraphics[width=1cm]{logo}};
\end{tikzpicture}%
}
\chead{\vspace{\fill}\Title\vspace{\fill}}
\rhead{\vspace{\fill}\Version\vspace{\fill}}
%quad è una spaziatura
%---------------------------------------%

\begin{table}[!h]
	\caption{Versioni}
	\footnotesize
	\begin{center}
		\begin{tabular}{ l l l l p{6cm} }
			\hline                                                                \\[-2ex]
			Ver. & Data       & Redattore        & Verificatore & Descrizione     \\
			\\[-2ex] \hline \\[-1.5ex]
			1.0  & 2024-06-08 & Davide Malgarise & Elena Ferro  & Stesura verbale \\
			\\[-1.5ex] \hline
		\end{tabular}
	\end{center}
\end{table}

\newpage

\tableofcontents

\newpage

\section{Dettagli della riunione}


\textbf{Sede della riunione}: Google Meet\\
\textbf{Orario di inizio}: 16:00\\
\textbf{Orario di fine}: 16:45\\

\begin{flushleft}
	\begin{table}[!h]
		\begin{tabular}{ |l|l|l| }
			\hline
			\textbf{Partecipante} & \textbf{Ruolo} & \textbf{Presenza} \\
			\hline
			Antonio Benetazzo     &                & Presente          \\
			Davide Malgarise      & Redattore      & Presente          \\
			Elena Ferro           & Verificatore   & Presente          \\
			Leonardo Baldo        &                & Presente          \\
			Matteo Tiozzo         &                & Presente          \\
			Raul Seganfreddo      &                & Presente          \\
			Valerio Occhinegro    & Amministratore & Presente          \\
			\hline
		\end{tabular}
	\end{table}
	\textbf{Partecipanti esterni}: Andrea Dorigo, Daniele Zorzi.\\
\end{flushleft}

\section{Ordine del giorno}
\begin{itemize}
	\item Stato di sviluppo del prodotto;
	\item Discussione sull'architettura del prodotto;
	\item Obiettivi per il settimo sprint;
	\item Decisioni prese e conclusioni.
\end{itemize}

\newpage

\section{Verbale}

\subsection{Stato di sviluppo del prodotto}
Come prima cosa, il responsabile del periodo appena concluso, Elena Ferro, riferisce al proponente l'esito della prima parte di valutazione RTB,
avvenuta in data 29 maggio in presenza del professore Riccardo Cardin. Il gruppo ha ricevuto un feedback positivo e può procedere con la seconda
parte di valutazione anche se sono state richieste alcune migliorie. Di conseguenza, assieme al proponente, il team discute e analizza
i requisiti funzionali, qualitativi e di vincolo. Ne risulta che il grado di specificità è adeguato e soddisfacente, anche se è consigliata l'aggiunta
di requisiti di performance e tempi di risposta.

\subsection{Discussione sull'architettura del prodotto}
L'incontro prosegue con la presentazione dell'architettura, che viene valutata positivamente e in linea con il progetto.
Quindi il gruppo espone quindi due dubbi emersi durante la fase di progettazione.\\\linebreak
Il primo riguarda la gestione dei dati provenienti dai sensori: nell'architettura progettata, i dati grezzi pubblicati su Redpanda vengono consumati contemporaneamente da \textit{Apache Flink}, che effettua aggregazioni,
e dal \textit{sink connector} per \textit{Clickhouse}, che salva i dati in apposite tabelle. Il proponente conferma la correttezza di questa scelta.\\\linebreak
Il secondo dubbio riguarda le operazioni da eseguire con \textit{Apache Flink}. Il proponente suggerisce di utilizzare questo strumento solo per aggregazioni complesse, come KPI o calcoli che coinvolgono dati provenienti da sensori diversi.
Di conseguenza, per computazioni più semplici, come la media per un certo periodo di tempo, è consigliato utilizzare le \textit{materialized view} di \textit{Clickhouse}, come già svolto fino a questo momento.

\subsection{Obiettivi per il settimo sprint}
Sono stati definiti i seguenti obiettivi per il settimo sprint:
\begin{itemize}
	\item Implementare due nuovi sensori;
	\item Iniziare la fase di test;
	\item Effettuare prove di carico e registrare i risultati;
	\item Aggiungere requisiti di performance.
\end{itemize}

\subsection{Decisioni prese e conclusioni}
In accordo con il proponente, si mantiene la durata del periodo di lavoro di una settimana e dunque si fissa il prossimo incontro in data
12 giugno 2024. La riunione si conclude con una breve discussione sulle possibile tempistiche di completamento e consegna del progetto,
auspicabile per la prima metà di luglio.

\begin{table}[b]
	\begin{tabular}{@{}p{5cm}p{10cm}@{}}
		Data:                 & \hrulefill \\
		                      &            \\
		                      &            \\
		Firma del proponente: & \hrulefill \\
	\end{tabular}
\end{table}

\end{document}
